\documentclass[conference]{IEEEtran}
\IEEEoverridecommandlockouts

\usepackage{cite}
\usepackage{amsmath,amssymb,amsfonts}
\usepackage{graphicx}
\usepackage{textcomp}
\usepackage{xcolor}
\usepackage{booktabs}
\usepackage{multirow}
\usepackage{microtype}

\def\BibTeX{{\rm B\kern-.05em{\sc i\kern-.025em b}\kern-.08em
    T\kern-.1667em\lower.7ex\hbox{E}\kern-.125emX}}

\begin{document}

\title{MinImageScorer: Error-Driven Difficulty Scoring\\
for Guided Copy-Paste Augmentation in\\
Agricultural Object Detection}

% TODO: Fill in your name, supervisor name, affiliation
\author{
\IEEEauthorblockN{Author Name\textsuperscript{1},
                  Supervisor Name\textsuperscript{1}}
\IEEEauthorblockA{
\textsuperscript{1}Department of XXX, University of XXX, Country\\
email@university.edu}
}

\maketitle

% ----------------------------------------------------------
\begin{abstract}
Copy-paste augmentation is widely applied to improve
object detection but is rarely evaluated against the
specific failure mode a dataset exhibits. We propose
\textbf{MinImageScorer}, an error-driven framework that
scores per-object detection difficulty from baseline model
failures using a cost function combining prediction
confidence and spatial accuracy, with weights
auto-calibrated via Pearson correlation optimization.
We validate the score on two agricultural benchmarks ---
MinneApple and Global Wheat Head (GWHD) --- and use it to
guide copy-paste augmentation. Experiments across three
random seeds show that score-guided selection outperforms
random selection on MinneApple ($\Delta$AP\,=\,+0.013)
but neither condition recovers the no-augmentation
baseline. Feature-space analysis confirms augmented images
shift the training distribution away from the test
distribution, explaining the consistent degradation.
On GWHD, domain shift dominates and copy-paste provides
no benefit under either selection strategy. Our findings
establish that the utility of copy-paste augmentation
depends on the dataset's underlying failure mode, and
that MinImageScorer correctly identifies the objects
responsible for that failure.
\end{abstract}

\begin{IEEEkeywords}
object detection, copy-paste augmentation, difficulty
scoring, agricultural detection, small objects,
domain shift
\end{IEEEkeywords}

% ----------------------------------------------------------
\section{Introduction}

Agricultural object detection benchmarks such as
MinneApple~\cite{hani2020minneapple} and Global Wheat
Head (GWHD)~\cite{david2021global} present distinct
structural failure modes. MinneApple is dominated by
high-density small objects near the resolution floor of
standard backbones, while GWHD contains images from
multiple geographic sources with significant cross-domain
visual variation. These failure modes differ fundamentally
from the object \textit{underrepresentation} problem that
Kisantal~\textit{et al.}~\cite{kisantal2019augmentation}
addressed on MS COCO, where copy-paste augmentation
produced measurable gains by increasing small object
anchor matches during training.

A critical gap remains: practitioners routinely apply
copy-paste augmentation without first diagnosing whether
the dataset's failure mode is addressable by adding more
object instances. When the failure is structural ---
resolution-limited detectability or cross-domain shift ---
copy-paste may corrupt the training distribution rather
than improve it. No prior work provides an empirical
framework to diagnose this before augmentation is applied.

We make the following contributions:
\begin{enumerate}
    \item \textbf{MinImageScorer}: a per-object difficulty
    score derived from model failure signals, with weights
    auto-calibrated via Pearson correlation optimization,
    requiring no domain-specific assumptions.
    \item \textbf{Empirical validation} that the score
    captures dataset-specific failure modes on both
    MinneApple and GWHD.
    \item \textbf{Controlled augmentation investigation}
    comparing score-guided and random copy-paste across
    three seeds, with feature-space analysis explaining
    the observed performance patterns.
\end{enumerate}

% ----------------------------------------------------------
\section{Methodology}

\subsection{MinImageScorer: Formulation}

Let $G(I)$ be the ground-truth objects in image $I$
and $\mathcal{P}$ the model predictions. For object $g$,
define $\mathcal{P}(g;\tau)$ as predictions with
$\mathrm{IoU}(p,g)\ge\tau$. The \textit{object difficulty
score} is:

\begin{equation}
S_{\mathrm{obj}}(g)=
\begin{cases}
\displaystyle\min_{p\in\mathcal{P}(g;\tau)}
\!\bigl[\alpha(1{-}\mathrm{conf}(p))
      +\beta(1{-}\mathrm{IoU}(p,g))\bigr],
& \mathcal{P}(g;\tau)\!\neq\!\varnothing,\\
\alpha+\beta, & \mathcal{P}(g;\tau)=\varnothing.
\end{cases}
\label{eq:obj_score}
\end{equation}

Missed detections receive the maximum score
$\alpha{+}\beta$. The \textit{image-level score}
aggregates object difficulties and penalises false
positives:

\begin{equation}
S_{\mathrm{img}}(I)=
w_1\!\cdot\!\frac{1}{|G(I)|}\!\sum_{g\in G(I)}\!
S_{\mathrm{obj}}(g)
\;+\;w_2\!\cdot\!\mathrm{FPrate}(I;\gamma,\tau).
\label{eq:img_score}
\end{equation}

Weights $\alpha,\beta,w_1,w_2$ are calibrated by
maximizing the Pearson correlation between score
distributions and per-image miss rate and false positive
rate on the training set.

\subsection{Score-Guided Copy-Paste Pipeline}

The pipeline proceeds as: (1) train a baseline
YOLOv8s~\cite{yolov8} on the original training set with
all built-in augmentations disabled; (2) run inference
on the training set to collect confidence and IoU signals;
(3) compute $S_{\mathrm{obj}}$ and $S_{\mathrm{img}}$
for all objects and images; (4) sample objects
proportionally to their scores using exponential weighting
($\alpha{=}2$) for score-guided, or uniformly for random;
(5) paste 2--8 objects per augmented image within the
source image (same-image-only); (6) retrain from COCO
pretrained weights on the augmented set.
Both conditions are identical in all parameters except
the selection policy, ensuring a controlled comparison.

% ----------------------------------------------------------
\section{Experiments}

\subsection{Setup}

\textbf{Datasets.}
MinneApple contains
% TODO: insert train/val/test image counts and
% total object instance counts here
images of apple trees with dense, small-object
annotations. GWHD contains
% TODO: insert GWHD train/val/test counts
images of wheat fields collected across multiple
geographic sources (Australia, Japan, India, and others).

\textbf{Model and training.} All models use YOLOv8s
initialized from COCO pretrained weights. Training
settings: MinneApple 200 epochs / patience 10; GWHD 100
epochs / patience 8; batch size 16; input size
1024$\times$1024; learning rate $10^{-2}$.
All Ultralytics built-in augmentations (mosaic, mixup,
HSV jitter, flips) are set to zero to isolate the
effect of copy-paste.

\textbf{Augmentation.} Dataset ratio 0.3; scale jitter
$[0.9,1.1]$ (MinneApple), $[0.8,1.2]$ (GWHD); no
rotation; overlap avoidance with 5\,px margin; no
blending. Results are reported as means over seeds
$\{123,456,789\}$.

\subsection{Results}

Table~\ref{tab:results} reports 3-seed mean COCO AP
for both datasets.

\begin{table}[t]
\centering
\caption{3-seed mean COCO AP on original test sets.
Bold = best copy-paste condition per dataset.}
\label{tab:results}
\resizebox{\columnwidth}{!}{%
\begin{tabular}{@{}llrrrrr@{}}
\toprule
\textbf{Dataset} & \textbf{Condition}
  & \textbf{AP} & \textbf{AP\textsubscript{50}}
  & \textbf{AP\textsubscript{75}}
  & \textbf{AP\textsubscript{S}}
  & \textbf{AP\textsubscript{M}} \\
\midrule
\multirow{3}{*}{MinneApple}
  & Baseline  & 0.439 & 0.771 & 0.452 & 0.300 & 0.592 \\
  & Random CP & 0.415 & 0.741 & 0.427 & 0.279 & 0.563 \\
  & \textbf{Score CP}
               & \textbf{0.428}
               & \textbf{0.759}
               & \textbf{0.440}
               & \textbf{0.286}
               & \textbf{0.581} \\
\midrule
\multirow{3}{*}{GWHD}
  & Baseline  & 0.507 & 0.922 & 0.491 & 0.136 & 0.501 \\
  & Random CP & 0.486 & 0.905 & 0.456 & 0.118 & 0.480 \\
  & \textbf{Score CP}
               & \textbf{0.484}
               & \textbf{0.905}
               & \textbf{0.455}
               & \textbf{0.098}
               & \textbf{0.477} \\
\bottomrule
\end{tabular}%
}
\end{table}

On MinneApple, score-guided copy-paste outperforms random
copy-paste by $\Delta\mathrm{AP}\,{=}\,{+}0.013$,
with consistent gains across AP, AP$_{50}$, AP$_{75}$,
AP$_S$, and AP$_M$. This indicates the score successfully
directs augmentation toward harder, smaller objects that
provide greater gradient signal. However, both copy-paste
conditions remain below baseline ($-0.024$ for Score CP),
showing that copy-paste adds distributional noise that
outweighs its selective benefit.

On GWHD, score-guided and random selection are
statistically indistinguishable (gap: $-0.002$ AP),
and both degrade below baseline ($-0.023$). The failure
pattern differs from MinneApple: score-guided sampling
concentrates pasted objects from a narrow subset of
source domains (predominantly \textit{arvalis}), further
amplifying the domain mismatch already present in the
dataset rather than addressing it.

\subsection{Feature-Space Analysis}

% TODO: Insert your feature space figure here.
% Suggested caption below — edit to match your actual figure.
\begin{figure}[t]
\centering
% \includegraphics[width=\columnwidth]{feature_space_combined.png}
\framebox{\parbox{0.9\columnwidth}{\centering
\vspace{1.8cm}
\small FIGURE PLACEHOLDER:\\
Image-level UMAP projections for MinneApple (left)\\
and GWHD (right). Groups: original train, augmented\\
(random CP), augmented (score CP), and test images.\\
\vspace{1.8cm}}}
\caption{Image-level backbone feature projections (UMAP,
P4 layer). Augmented images occupy regions that are
systematically farther from the test distribution than
original training images, confirming distributional
corruption under both augmentation strategies.
Score-guided augmented images are displaced further on
GWHD due to domain-minority object concentration.}
\label{fig:feature_space}
\end{figure}

Fig.~\ref{fig:feature_space} shows image-level backbone
embeddings (YOLOv8s P4, global average pool, UMAP
projection) for four groups: original training images,
augmented images under random CP, augmented images under
score CP, and test images. On both datasets, augmented
images occupy a region of feature space that is displaced
from the test distribution relative to the original
training images.

% TODO: Insert centroid distance values from your analysis
% Example: "The centroid distance from augmented score CP
% to the test set was X.XX vs X.XX for original training
% images on MinneApple."
Quantitatively, the centroid distance from augmented
images to the test distribution exceeds that of original
training images on both datasets
(MinneApple: \textit{[value]} vs \textit{[value]};
GWHD: \textit{[value]} vs \textit{[value]}),
directly explaining the consistent performance
degradation. On GWHD, score-guided augmented images show
greater displacement than random augmented images,
consistent with domain-minority object concentration.

% ----------------------------------------------------------
\section{Discussion}

\textbf{Why copy-paste fails on these datasets.}
Kisantal~\textit{et al.}~\cite{kisantal2019augmentation}
demonstrated that copy-paste improves AP$_S$ on MS COCO
by increasing anchor matches for underrepresented small
objects. This benefit presupposes that: (1) hard objects
are above the backbone feature resolution floor, and
(2) pasted objects come from the same visual domain as
the target scene. Neither condition holds reliably on
MinneApple or GWHD. On MinneApple, the hardest objects
identified by MinImageScorer are predominantly
near the resolution limit of the backbone, making them
weakly represented in feature maps regardless of paste
frequency. On GWHD, even within-image pasting fails
because the feature-space mismatch originates at the
image level, not the object level.

\textbf{Diagnostic value of the score.}
Despite producing no absolute improvement, MinImageScorer
demonstrates clear utility: it correctly ranks objects
by detectability (score correlates with miss rate;
% TODO: insert Pearson r values
$r\,{=}\,[\text{value}]$ on MinneApple,
$r\,{=}\,[\text{value}]$ on GWHD),
and on MinneApple it selects objects that produce
a +0.013 AP relative gain over uniform random selection.
The score identifies the right objects; the limitation is
that copy-paste cannot make those objects more learnable
when the failure mode is structural.

\textbf{Established boundary condition.}
Score-guided copy-paste provides a positive relative
benefit when the selected objects are above the
resolution floor and within-distribution with the target
scene. It provides no benefit when the failure mode is
domain shift, and it amplifies the degradation on
AP$_S$ when objects selected are predominantly
sub-resolution. This constitutes a practical
diagnostic criterion: run MinImageScorer, inspect the
size and domain distribution of high-score objects
before applying copy-paste augmentation.

% ----------------------------------------------------------
\section{Conclusion}

We proposed MinImageScorer, a model-failure-driven
difficulty scoring framework validated on two
agricultural detection benchmarks. Controlled experiments
show that score-guided copy-paste outperforms random
copy-paste on MinneApple ($+$0.013 AP) but that both
strategies degrade below baseline, with the degradation
confirmed by feature-space analysis to originate from
distributional shift between augmented training images
and the test set. On GWHD, domain shift makes copy-paste
ineffective under any selection strategy.
The practical implication is that copy-paste augmentation
should be preceded by failure mode diagnosis: it is
effective for underrepresentation (COCO-like datasets)
but not for resolution-limited or domain-shifted datasets.
Future work includes online augmentation to reduce
distributional rigidity and domain-constrained pasting
for multi-source benchmarks.

% ----------------------------------------------------------
\begin{thebibliography}{00}

\bibitem{hani2020minneapple}
N.~Hani, P.~Roy, and B.~Bhanu,
``MinneApple: A Benchmark Dataset for Apple Detection
and Segmentation,''
\textit{IEEE Robotics and Automation Letters},
vol.~5, no.~2, pp.~852--858, 2020.

\bibitem{david2021global}
E.~David \textit{et al.},
``Global Wheat Head Detection (GWHD) Dataset,''
\textit{Plant Phenomics}, vol.~2021, 2021.

\bibitem{kisantal2019augmentation}
M.~Kisantal \textit{et al.},
``Augmentation for Small Object Detection,''
in \textit{Proc. 9th Int. Conf. Advances in Computing,
Communication and Informatics}, 2019.

\bibitem{yolov8}
G.~Jocher \textit{et al.},
``Ultralytics YOLOv8,''
GitHub, 2023. [Online].
Available: https://github.com/ultralytics/ultralytics

% TODO: Add additional references as needed
% Suggested additions:
% - Global Wheat Head challenge paper (Rémi Dufumier et al.)
% - Simple Copy-Paste (Ghiasi et al., CVPR 2021)
% - Focal Loss / RetinaNet (Lin et al.) if cited
% - OHEM (Shrivastava et al.) if cited

\end{thebibliography}

\end{document}